\documentclass[a4paper,11pt]{article}

\usepackage{fontspec}
\setmainfont{Carlito}
\usepackage[margin=0.6in, top=0.5in, bottom=0.5in]{geometry}
\usepackage{titlesec}
\usepackage{enumitem}
\usepackage{hyperref}
\usepackage{xcolor}
\usepackage{microtype}
\usepackage{needspace}

\hyphenpenalty=10000 \exhyphenpenalty=10000 \emergencystretch=3em

\definecolor{headernavy}{RGB}{12, 83, 166}
\definecolor{linkblue}{RGB}{0, 102, 204}
\definecolor{muted}{RGB}{96, 102, 110}

\hypersetup{colorlinks=true, linkcolor=linkblue, urlcolor=linkblue}

\titleformat{\section}
  {\color{headernavy}\normalsize\bfseries}
  {}{0em}{\MakeUppercase}[\color{headernavy}\titlerule]

\titlespacing{\section}{0pt}{8pt}{4pt}

\setlist[itemize]{noitemsep, topsep=1pt, leftmargin=1.3em, partopsep=0pt, parsep=1pt}

\pagestyle{empty}
\setlength{\parindent}{0pt}
\linespread{1.02}

% ---- building blocks: bold = organisation/project, italic = role/stack, grey = dates/places ----
% v3.5.3 (10 Nov 2026): employer (and place) on line 1, the title then the dates on line 2, so that `pdftotext -raw` reads every role as
% employer -> title -> dates -> bullets (ats_structure.role_order_problems). The dates were on line 1 before, read before the title.
\newcommand{\role}[4]{% org, location, dates, title
  \textbf{#1}{\color{muted}\,\textperiodcentered\,\textit{#2}}\\
  \textit{#4}\hfill{\color{muted}#3}\par}
\newcommand{\roleNoLoc}[3]{\textbf{#1}\\ \textit{#3}\hfill{\color{muted}#2}\par}
\newcommand{\project}[3]{% name, stack, right-hand tag
  \textbf{#1}\hfill#3\\ {\color{muted}\small\textit{#2}}\par}
\newcommand{\gh}[1]{\href{#1}{\small GitHub}}
\newcommand{\badge}[1]{{\color{headernavy}\small\textbf{#1}}}
\newcommand{\skillgroup}[2]{{\color{headernavy}\textbf{#1}}\\ #2\par\vspace{3pt}}
\newcommand{\gap}{\vspace{5pt}}
\newcommand{\num}[1]{\textbf{#1}}   % result numbers only


\begin{document}

% ================= HEADER =================
\begin{center}
  {\LARGE\bfseries Thiruvel Andagurunathan Pandian} \\[3pt]
  {\color{headernavy}\small\textbf{Data Scientist | Audited \num{859,494} training rows for quality and duplicates}} \\[3pt]
  {\small
  Bristol, UK\nobreak\ \,\textbar\, \href{mailto:apthiruvel@gmail.com}{apthiruvel@gmail.com}\nobreak\ \,\textbar\, \href{tel:+447741016954}{+44 7741 016954}\nobreak\ \,\textbar\, \href{https://www.linkedin.com/in/thiruvel-a-p/}{https://www.linkedin.com/in/thiruvel-a-p/}\nobreak\ \,\textbar\, \href{https://github.com/Thiruvel-AP}{https://github.com/Thiruvel-AP}\nobreak\ \,\textbar\, \href{https://thiruvel-ap.github.io/PersonalPortfolio/}{https://thiruvel-ap.github.io/PersonalPortfolio/}}\par\vspace{4pt}
  \noindent{} Right to Work in UK \,\textbar\, Available to join Immediately
\end{center}

% ================= PROFILE =================
\needspace{5\baselineskip}\section{Profile}
Data scientist who checks the data before the model: audited 859,494 training examples for gaps and duplicates, and combined 14 research datasets for my AstraZeneca-affiliated MSc dissertation. Deploys on AWS (Lambda) with CI/CD and MLflow experiment tracking. Builds the model a problem needs and says plainly where it falls short: customer-churn prediction, price forecasting, text grouping, and models for 11,459 genes that scored 0.807 on training data but 0.545 on unseen data. Built an AI assistant on AWS for an MSc dissertation with an AstraZeneca affiliation; it checks its answers against official databases and is tested by about 65 automated checks.

% ================= SKILLS =================
\needspace{5\baselineskip}\section{Skills}
\skillgroup{Data analysis}{data analysis, eda, exploratory data analysis, data audit, data quality, Explored}
\skillgroup{Machine Learning, Deep Learning \& Statistics}{Machine learning concepts, statistical analysis (descriptive statistics, correlation analysis), predictive modelling (gradient boosting with XGBoost, random forest, scikit-learn), clustering (K-Means, HAC, DBSCAN, Leiden), deep learning (PyTorch, TensorFlow/Keras, LSTM, GRU, BiRNN, CNN), time-series forecasting (cyclic time encoding, TimeSeriesSplit), anomaly detection (Isolation Forest), class-imbalance handling (SMOTE), hyper-parameter tuning (grid search, two-stage Bayesian), model evaluation (held-out, cross-validation, Pearson, MCC), statistical inference and experimental design \textit{(confidence intervals, bootstrap, hypothesis testing; MSc coursework)}, MLflow experiment tracking}
\skillgroup{Python, Data Quality \& Cloud}{Python, pandas, NumPy, SciPy, Matplotlib, Seaborn, Plotly, SQL, PostgreSQL, R (tidymodels), identifier crosswalks, parquet, reproducible pipelines, AWS, Docker, Kubernetes (AWS EKS), serverless deployment (AWS Lambda), containerisation, GitHub Actions CI/CD, FastAPI, TensorFlow Serving, pytest, Git}
\skillgroup{NLP, LLMs, Embeddings \& Evaluation}{NLP (text preprocessing, TF-IDF, SBERT embeddings, text clustering), Hugging Face Transformers, large language models (Claude, OpenAI, open-weight via Ollama), prompt engineering, RAG with embeddings (ChromaDB), agents (LangChain, LangGraph), team LoRA fine-tune of Llama, LLM evaluation (hallucination-rate analysis, deterministic guardrails)}

% ================= EXPERIENCE =================
\needspace{5\baselineskip}\section{Work Experience}
\needspace{4\baselineskip}
\role{AstraZeneca}{Bristol, UK (Remote)}{May 2026 – Sep 2026}{Data Scientist, MSc Dissertation Project}
\begin{itemize}
  \item Needed an LLM a research team could trust before it shipped: built a tool-calling LLM agent (LangChain, AWS Bedrock, Claude Haiku 4.5) for a multi-omics cancer research platform, with three tools for gene alias resolution, scoring-methodology explanation and dataset provenance.
  \item Integrated the agent into a React front end with real-time gene alias enrichment.
  \item Found where the assistant went wrong with about 65 automated checks at three levels, fixed failures layer by layer, and reused the checks after changing AI provider (pytest).
  \item Combined 14 separate research datasets so the assistant only used trusted data, then put it online on AWS and documented how it works (Python, AWS Lambda).
\end{itemize}
\gap
\needspace{4\baselineskip}
\role{180 Degrees Consulting Bristol}{Bristol, UK (Remote)}{Sep 2025 – Dec 2025}{Data Consultant (AI Engineering \& AI System)}
\begin{itemize}
  \item Led the design in a five-person team that turned a startup client\textquotesingle{}s vague idea into a working AI prototype in three months; senior leadership rated it 5 stars (Python, MCP).
  \item Proposed the request flow: a user\textquotesingle{}s text goes through the API and a customised prompt to the MCP server, where one LLM instance writes GraphQL queries against the client\textquotesingle{}s Convex database and a second instance runs the statistical analysis before the result returns to the user.
  \item Helped the client\textquotesingle{}s engineers put the tool into everyday use and reviewed their code changes with the CTO; it is \num{now in live use} (TypeScript, OpenAI).
\end{itemize}
\gap
\needspace{4\baselineskip}
\role{Freelance}{Chennai, India}{Jan 2025 – Aug 2025}{AI/ML Engineer}
\begin{itemize}
  \item Working independently for a private client, built agentic RAG workflows (LangChain, LangGraph, ChromaDB) self-hosted on the client\textquotesingle{}s own server, NLP pipelines (Hugging Face Transformers, Ollama) and scikit-learn pipelines, served ML models as REST APIs (FastAPI, Docker), and contributed to a team LoRA fine-tune of a Llama model.
  \item Built a private client an AI finance assistant that reads his income, spending, loans and investments and suggests next steps, self-hosted to keep costs low (Llama, Ollama).
  \item Designed reproducible scikit-learn pipelines with feature engineering (temporal features, interaction terms, target encoding), model selection, cross-validation and hyperparameter tuning across classification and regression tasks.
\end{itemize}
\gap
\needspace{4\baselineskip}
\role{Avasoft}{Chennai, India}{Jan 2023 – Sep 2024}{Software Development Engineer Trainee, then Software Development Engineer}
\begin{itemize}
  \item Designed Core Data schemas for user profiles, transaction histories and event bookings with CRUD and session-state management.
  \item Helped build a money-management app for a company\textquotesingle{}s 10,000+ employees, linking live stock prices and bank accounts to a backend of 200+ endpoints; shipped the first release (Flutter, AWS).
  \item \num{Shipped features in a production fintech (financial-services) application} as part of a team: integrated Twelve Data (real-time market feeds) and Plaid (bank linking, transaction imports) via REST APIs, with SSL pinning and AWS Cognito SSO.
\end{itemize}

% ================= PROJECTS =================
\needspace{5\baselineskip}\section{Projects}
\needspace{4\baselineskip}
\setbox0=\hbox{\textbf{llama3-8b-qlora-math-t4: QLoRA fine-tuning of Llama 3.1 8B for maths under a token budget}\hspace{1em}\url{https://github.com/Thiruvel-AP/llama3-8b-qlora-math-t4}}%
\setbox2=\hbox{{\small\textit{Python, Unsloth, PEFT, TRL, Hugging Face, PyTorch, sentence-transformers, scikit-learn}}\hspace{1em}{\small\url{https://github.com/Thiruvel-AP/llama3-8b-qlora-math-t4}}}%
\ifdim\wd0>\linewidth
\ifdim\wd2>\linewidth
\textbf{llama3-8b-qlora-math-t4: QLoRA fine-tuning of Llama 3.1 8B for maths under a token budget}\\ {\color{muted}\small\textit{Python, Unsloth, PEFT, TRL, Hugging Face, PyTorch, sentence-transformers, scikit-learn}}\\ \hspace*{\fill}{\small\url{https://github.com/Thiruvel-AP/llama3-8b-qlora-math-t4}}\par
\else
\textbf{llama3-8b-qlora-math-t4: QLoRA fine-tuning of Llama 3.1 8B for maths under a token budget}\\ {\color{muted}\small\textit{Python, Unsloth, PEFT, TRL, Hugging Face, PyTorch, sentence-transformers, scikit-learn}}\hfill{\small\url{https://github.com/Thiruvel-AP/llama3-8b-qlora-math-t4}}\par
\fi
\else
\project{llama3-8b-qlora-math-t4: QLoRA fine-tuning of Llama 3.1 8B for maths under a token budget}{Python, Unsloth, PEFT, TRL, Hugging Face, PyTorch, sentence-transformers, scikit-learn}{\url{https://github.com/Thiruvel-AP/llama3-8b-qlora-math-t4}}
\fi
\begin{itemize}
  \item Checked all 859,494 training examples for gaps and copies, found 44,369 extra duplicates and hand-reviewed 181 suspect records before training (Python).
\end{itemize}
\gap
\needspace{4\baselineskip}
\setbox0=\hbox{\textbf{Mining Insights from Customer Feedback: NLP \& unsupervised analytics (MSc team project)}\hspace{1em}\url{https://github.com/Thiruvel-AP/TextAnalytics}}%
\setbox2=\hbox{{\small\textit{Python, scikit-learn, SBERT, K-Means, HAC, DBSCAN}}\hspace{1em}{\small\url{https://github.com/Thiruvel-AP/TextAnalytics}}}%
\ifdim\wd0>\linewidth
\ifdim\wd2>\linewidth
\textbf{Mining Insights from Customer Feedback: NLP \& unsupervised analytics (MSc team project)}\\ {\color{muted}\small\textit{Python, scikit-learn, SBERT, K-Means, HAC, DBSCAN}}\\ \hspace*{\fill}{\small\url{https://github.com/Thiruvel-AP/TextAnalytics}}\par
\else
\textbf{Mining Insights from Customer Feedback: NLP \& unsupervised analytics (MSc team project)}\\ {\color{muted}\small\textit{Python, scikit-learn, SBERT, K-Means, HAC, DBSCAN}}\hfill{\small\url{https://github.com/Thiruvel-AP/TextAnalytics}}\par
\fi
\else
\project{Mining Insights from Customer Feedback: NLP \& unsupervised analytics (MSc team project)}{Python, scikit-learn, SBERT, K-Means, HAC, DBSCAN}{\url{https://github.com/Thiruvel-AP/TextAnalytics}}
\fi
\begin{itemize}
  \item Needed structure in unstructured airline support tickets: built the text-cleaning and grouping steps for airline support tickets, and traced weak groupings (scores below 0.30) to repeated template wording (Python, scikit-learn).
\end{itemize}
\gap
\needspace{4\baselineskip}
\setbox0=\hbox{\textbf{CoverFit: insurance benchmarking API (Bristol × UWE AI Hackathon)}\hspace{1em}\url{https://github.com/Thiruvel-AP/AIHackathon-Phoenix}}%
\setbox2=\hbox{{\small\textit{FastAPI, pandas, LiteLLM, Pydantic, pytest}}\hspace{1em}{\small\url{https://github.com/Thiruvel-AP/AIHackathon-Phoenix}}}%
\ifdim\wd0>\linewidth
\ifdim\wd2>\linewidth
\textbf{CoverFit: insurance benchmarking API (Bristol × UWE AI Hackathon)}\\ {\color{muted}\small\textit{FastAPI, pandas, LiteLLM, Pydantic, pytest}}\\ \hspace*{\fill}{\small\url{https://github.com/Thiruvel-AP/AIHackathon-Phoenix}}\par
\else
\textbf{CoverFit: insurance benchmarking API (Bristol × UWE AI Hackathon)}\\ {\color{muted}\small\textit{FastAPI, pandas, LiteLLM, Pydantic, pytest}}\hfill{\small\url{https://github.com/Thiruvel-AP/AIHackathon-Phoenix}}\par
\fi
\else
\project{CoverFit: insurance benchmarking API (Bristol × UWE AI Hackathon)}{FastAPI, pandas, LiteLLM, Pydantic, pytest}{\url{https://github.com/Thiruvel-AP/AIHackathon-Phoenix}}
\fi
\begin{itemize}
  \item For an insurance hackathon brief, built a tool that finds comparable companies and explains the match in under 3 seconds for about £0.003 per query (FastAPI, Gemini).
\end{itemize}

% ================= EDUCATION =================
\needspace{5\baselineskip}\section{Education}
\noindent\textit{MSc Data Science}, \textbf{University of Bristol, UK}\unskip\nobreak\hfill\penalty50\hskip1em\null\nobreak\hfill\mbox{Sep 2025 – Sep 2026}\par
\noindent{\small Dissertation with AstraZeneca affiliation (GeneTraceAI)}\par
\noindent\textit{B.E. Electrical \& Electronics Engineering}, \textbf{St. Joseph\textquotesingle{}s College of Engineering}, First Class with Distinction\unskip\nobreak\hfill\penalty50\hskip1em\null\nobreak\hfill\mbox{Aug 2019 – Apr 2023}\par

% ================= CERTIFICATIONS & ACHIEVEMENTS =================
\needspace{5\baselineskip}\section{Certifications \& Achievements}
\begin{itemize}
  \item Model Context Protocol: Advanced Topics, Anthropic, \textit{Jun 2026}
  \item Introduction to Model Context Protocol, Anthropic, \textit{Jun 2026}
  \item AWS Academy Cloud Foundations, \textit{Nov 2025}
  \item 180 Degrees Consulting Data Science Consultant Certificate, \textit{Dec 2025}
  \item Bristol PLUS Award, \textit{Mar 2026}
\end{itemize}

\end{document}